\bmtchapitre{Barycentres et lignes de niveau}{Construire et calculer un barycentre, regrouper des points pondérés et déterminer des lieux géométriques.}{Géométrie}
\label{t11s-c14}
Dans ce chapitre, les calculs de longueurs, d’angles et d’aires à partir de coordonnées se font dans un repère orthonormé.
\begin{revision}Chasles, produit scalaire, normes, milieux et équations de cercles.\end{revision}
\begin{automatismes}\exo\label{t11s-c14-auto}Calculer le milieu de $A(0;1)$ et $B(6;3)$. Développer $(x-2)^2$.\end{automatismes}
\begin{corrige}[exercice~\ref{t11s-c14-auto}]\label{sol-t11s-c14-auto}Milieu $(3;2)$; $x^2-4x+4$.\end{corrige}

\section{Deux, trois ou quatre points pondérés}
\begin{definition}Pour des points $A_i$ et des poids réels $\alpha_i$ dont la somme $s$ est non nulle, le barycentre $G$ est l’unique point tel que $\sum_i\alpha_i\vect{GA_i}=\vec0$.\end{definition}
\begin{propriete}[existence, coordonnées et réduction]
Pour tout point $O$, $\vect{OG}=\dfrac1s\sum_i\alpha_i\vect{OA_i}$. Pour tout $M$,
\[
\sum_i\alpha_i\vect{MA_i}=s\vect{MG}.
\]
Les coordonnées de $G$ sont les moyennes pondérées des coordonnées des points. Multiplier tous les poids par un même réel non nul ne change pas $G$.
\end{propriete}
\begin{demonstration}Par Chasles, $\vect{GA_i}=\vect{GO}+\vect{OA_i}$. L’équation du barycentre impose $s\vect{OG}=\sum\alpha_i\vect{OA_i}$ : un unique vecteur définit un unique point. La réduction résulte de $\vect{MA_i}=\vect{MG}+\vect{GA_i}$.\end{demonstration}
\begin{exemple}Pour $(A,1),(B,2)$, $\vect{AG}=\tfrac23\vect{AB}$. Avec des poids $1,-2$, la somme vaut $-1$ et $\vect{AG}=2\vect{AB}$ : le barycentre peut être à l’extérieur du segment. Si les poids sont $1,-1$ et $A\ne B$, aucun point ne satisfait la définition, car la somme vectorielle est $\vect{BA}\ne\vec0$.\end{exemple}
\begin{propriete}[associativité]Un groupe de points dont la somme des poids est non nulle peut être remplacé par son barycentre affecté de cette somme. La somme totale doit rester non nulle.\end{propriete}
\begin{demonstration}La réduction de ce groupe remplace exactement sa contribution dans la somme des vecteurs. L’équation finale ne change donc pas.\end{demonstration}
\section{Sommes de distances au carré}
\begin{propriete}[identité de réduction]Pour le barycentre $G$ et $s\ne0$,
\[
\sum_i\alpha_i MA_i^2=sMG^2+\sum_i\alpha_i GA_i^2.
\]
Les poids peuvent être négatifs; on ne doit alors pas supposer que la somme est positive.
\end{propriete}
\begin{demonstration}Développer le carré de $\vect{MA_i}=\vect{MG}+\vect{GA_i}$, multiplier par $\alpha_i$ et sommer. Le terme croisé vaut $2\vect{MG}\cdot\sum\alpha_i\vect{GA_i}=0$.\end{demonstration}
\section{Reconnaître les lieux}
\begin{propriete}[lignes de niveau du plan]
$MA=k$ donne un cercle si $k>0$, le point $A$ si $k=0$, l’ensemble vide si $k<0$.
Pour $\vecv\ne\vec0$, $\vect{MA}\cdot\vecv=0$ est la droite par $A$ perpendiculaire à $\vecv$; pour $\vecv=\vec0$, tout le plan.
Si $I$ est milieu de $[AB]$,
\[
\vect{MA}\cdot\vect{MB}=MI^2-\frac{AB^2}{4}.
\]
Le lieu $\vect{MA}\cdot\vect{MB}=k$ est donc un cercle, un point ou vide selon le signe de $k+AB^2/4$.
\end{propriete}
\begin{demonstration}Écrire $\vect{MA}=\vect{MI}+\vect{IA}$ et $\vect{MB}=\vect{MI}-\vect{IA}$. Le produit vaut $MI^2-IA^2$, et $IA=AB/2$.\end{demonstration}
\begin{exemple}[rapport de distances]
Supposons $A\ne B$ et cherchons $MA=kMB$. Pour $k<0$, aucun point; pour $k=0$, le point $A$; pour $k=1$, la médiatrice de $[AB]$. Pour $k>0$, $k\ne1$, mettre au carré est équivalent; le lieu est un cercle d’Apollonius. Son centre est le barycentre de $(A,1),(B,-k^2)$ et son rayon vaut $kAB/|1-k^2|$.
Pour $A(0;0),B(3;0),k=2$, $x^2+y^2=4((x-3)^2+y^2)$ donne $(x-4)^2+y^2=4$ : centre $(4;0)$ et rayon $2$.
\end{exemple}

% ENRICHISSEMENT C14

\subsection*{Réduire une expression avant de nommer un lieu}
\begin{methode}{une équation de niveau en trois étapes}
Calculer la somme des poids et le barycentre si elle est non nulle. Réduire ensuite l’expression à $sMG^2+c$. Résoudre l’équation pour $MG^2$; une valeur positive donne un cercle, zéro un point, négative l’ensemble vide. Si $s<0$, le sens d’une inégalité serait renversé en divisant; ne pas supposer que l’expression est une somme positive.
\end{methode}
\begin{exemple}[une somme pondérée et son minimum]
Avec $A(0;0)$, $B(6;0)$ et poids $2,1$, le barycentre est $G(2;0)$. La réduction donne
\[
2MA^2+MB^2=3MG^2+2GA^2+GB^2=3MG^2+24.
\]
Le minimum est $24$, atteint uniquement en $G$. Le niveau $12$ est vide; le niveau $24$ est $\{G\}$; le niveau $51$ est le cercle de centre $G$ et de rayon $3$. Il faut calculer la constante de réduction, pas seulement le centre.
\end{exemple}
\begin{methode}{construire en regroupant les poids}
Pour trois points, chercher un groupe dont le barycentre est facile à construire. Si ses poids ont une somme non nulle, remplacer ce groupe par son barycentre affecté de cette somme. Le barycentre final se construit alors sur une droite entre deux points. Vérifier les deux sommes avant les divisions.
\end{methode}
\begin{remarque}[somme nulle ne signifie pas expression inutilisable]
Quand les poids totalisent zéro, la formule du barycentre ne s’applique pas. L’expression initiale peut néanmoins se simplifier directement par Chasles ou en coordonnées. Son lieu n’est pas automatiquement vide; il faut étudier cette expression, sans inventer un centre.
\end{remarque}

\section{Exercices et problèmes}
\exo[Deux poids]\label{t11s-c14-e01}
Pour $A(1;2),B(7;-1)$, calculer le barycentre de $(A,2),(B,1)$.
\begin{corrige}[exercice~\ref{t11s-c14-e01}]\label{sol-t11s-c14-e01}
Somme $3\ne0$. $G=((2+7)/3;(4-1)/3)=(3;1)$.
\end{corrige}
\exo[Regrouper]\label{t11s-c14-e02}
Dans un triangle $ABC$, $I$ est milieu de $[BC]$. Montrer que le barycentre de $(A,2),(B,1),(C,1)$ est le milieu de $[AI]$.
\begin{corrige}[exercice~\ref{t11s-c14-e02}]\label{sol-t11s-c14-e02}
Remplacer $(B,1),(C,1)$ par $(I,2)$. Les deux poids $2,2$ donnent le milieu, somme totale $4$.
\end{corrige}
\exo[Quatre points]\label{t11s-c14-e03}
Pour le rectangle $A(0;0),B(4;0),C(4;2),D(0;2)$, trouver le barycentre des quatre points de poids $1$.
\begin{corrige}[exercice~\ref{t11s-c14-e03}]\label{sol-t11s-c14-e03}
Somme $4$; coordonnées $(2;1)$, centre du rectangle. Regrouper les deux diagonales donne le même résultat.
\end{corrige}
\exo[Lieu avec produit]\label{t11s-c14-e04}
Si $AB=6$, déterminer les lieux $\vect{MA}\cdot\vect{MB}=0$, puis $-9$, puis $-10$.
\begin{corrige}[exercice~\ref{t11s-c14-e04}]\label{sol-t11s-c14-e04}
Avec $I$ milieu, $MI^2=k+9$. Pour $0$, cercle de centre $I$, rayon $3$ (diamètre $AB$); pour $-9$, point $I$; pour $-10$, ensemble vide.
\end{corrige}
\exo[Somme de carrés]\label{t11s-c14-e05}
Pour $AB=4$ et $I$ milieu, déterminer le lieu $MA^2+MB^2=18$.
\begin{corrige}[exercice~\ref{t11s-c14-e05}]\label{sol-t11s-c14-e05}
Théorème de la médiane : $MA^2+MB^2=2MI^2+8$. Donc $MI^2=5$ : cercle de centre $I$, rayon $\sqrt5$.
\end{corrige}
\exo[Apollonius]\label{t11s-c14-e06}
Pour $A(0;0),B(3;0)$, déterminer le lieu $MA=\tfrac12 MB$ et donner le centre par barycentre.
\begin{corrige}[exercice~\ref{t11s-c14-e06}]\label{sol-t11s-c14-e06}
$4(x^2+y^2)=(x-3)^2+y^2$, donc $(x+1)^2+y^2=4$. Centre $(-1;0)$, rayon $2$; barycentre de $(A,1),(B,-1/4)$, somme $3/4$.
\end{corrige}
\exo[Poids nuls au total]\label{t11s-c14-e07}
Pourquoi ne peut-on appliquer la formule du barycentre aux poids $1,-1$ ? Distinguer $A\ne B$ et $A=B$.
\begin{corrige}[exercice~\ref{t11s-c14-e07}]\label{sol-t11s-c14-e07}
Division par somme nulle interdite. Si $A\ne B$, la somme est $\vect{BA}\ne0$, aucun point. Si $A=B$, elle vaut zéro pour tout point : aucune unicité, donc pas de barycentre au sens défini.
\end{corrige}
\exo[Distance et projection]\label{t11s-c14-e08}
Dans un repère orthonormé, $A(1;2)$ et $\vecv=(2;-1)$. Déterminer les lieux $MA=k$ pour $k=2$, $0$, $-1$, puis le lieu $\vect{MA}\cdot\vecv=0$. Que devient ce dernier lieu si $\vecv=\vec0$ ?
\begin{corrige}[exercice~\ref{t11s-c14-e08}]\label{sol-t11s-c14-e08}
Pour $k=2$, cercle de centre $A$ et rayon $2$; pour $k=0$, le seul point $A$; pour $k=-1$, ensemble vide, car une distance est positive ou nulle. Pour $M(x;y)$, le produit vaut $2(1-x)-(2-y)=y-2x$, donc le lieu est la droite $y=2x$, passant par $A$ et perpendiculaire à $(2;-1)$. Si $\vecv=\vec0$, le produit vaut toujours zéro et le lieu est tout le plan.
\end{corrige}

\exo[Des poids négatifs donnent un maximum]\label{t11s-c14-p01}
Dans un repère orthonormé, $A(0;0)$, $B(4;0)$ et $F(M)=MA^2-2MB^2$. Calculer le barycentre des poids $1,-2$, puis montrer que $F(M)=32-MG^2$. En déduire le maximum de $F$ et déterminer les lieux $F(M)=20$, $32$, $33$.
\begin{corrige}[exercice~\ref{t11s-c14-p01}]\label{sol-t11s-c14-p01}
La somme des poids vaut $-1$, et $G=(8;0)$. La constante est $GA^2-2GB^2=64-32=32$, donc $F=32-MG^2$. Le maximum $32$ est atteint seulement en $G$. Pour le niveau $20$, $MG^2=12$ : cercle de rayon $2\sqrt3$; pour $32$, point $G$; pour $33$, $MG^2=-1$, ensemble vide. La fonction peut être négative et n’a pas de minimum, car $MG$ peut croître sans borne.
\end{corrige}
\exo[Une différence de carrés sans barycentre]\label{t11s-c14-p02}
Avec $A(0;0)$, $B(4;0)$ et $M(x;y)$ dans un repère orthonormé, calculer $MA^2-MB^2$. Pour tout réel $k$, déterminer le lieu $MA^2-MB^2=k$. Retrouver la médiatrice pour $k=0$ et expliquer pourquoi les poids $1,-1$ ne définissent pas de barycentre.
\begin{corrige}[exercice~\ref{t11s-c14-p02}]\label{sol-t11s-c14-p02}
La différence vaut $x^2+y^2-[(x-4)^2+y^2]=8x-16$. Le lieu est donc la droite $x=(k+16)/8$, pour chaque $k\in\R$. Pour $k=0$, on obtient $x=2$, médiatrice de $[AB]$. Les poids ont somme nulle et $A\ne B$, donc pas de barycentre; cette absence n’empêche pas les lieux d’exister.
\end{corrige}

\begin{jemeteste}Je vérifie la somme des poids avant de diviser et je classe un lieu par le signe de son rayon au carré.\end{jemeteste}
\begin{notepourlaclasse}Construire d’abord le barycentre de poids positifs, puis faire apparaître un poids négatif. Les cas cercle/point/vide doivent figurer dans les corrections, sans ajouter une figure fictive pour un lieu vide.\end{notepourlaclasse}
